\subsection{生成模}

设 \(\mathrm{A}\) 是一个环。

定义 2. —— 若每个 A-模 N 都由从 M 到 N 的 A-线性映射的像生成，则称 A-模 M 为生成模。

设 \(M\) 是一个左 A-模。我们把 \(M\) 的对偶记作 \(M^*\)，并把 \(M \times M^*\) 上的典范双线性型（II, §2, No.~3, 第 233 页）记作

\[
(x, x ^ {*}) \mapsto \langle x, x ^ {*} \rangle = x ^ {*} (x).
\]

用 \(\tau(\mathrm{M})\) 表示 A 中形如 \(\sum_{i=1}^{n}\langle x_i,x_i^*\rangle\) 的元素组成的集合，其中 \(x_1,\ldots ,x_n\) 是 M 的元素，\(x_1^*,\ldots ,x_n^*\) 是 \(\mathrm{M}^*\) 的元素。它是 A 的一个双边理想，称为 M 的迹理想。A-模 \(\mathrm{A}_s\) 的迹理想是 A。一族 A-模 \((\mathrm{M}_i)_{i\in \mathrm{I}}\) 的直和的迹理想是理想 \(\sum_{i\in \mathrm{I}}\tau(\mathrm{M}_i)\)。若 M 是投射 A-模，则由 II, §2, No.~6, 第 238 页的命题 12，我们有 \(\mathrm{M} = \tau(\mathrm{M})\mathrm{M}\)。

定理 1. —— 设 M 是一个左 A-模。下列性质等价：

\begin{enumerate}
\def\labelenumi{(\roman{enumi})}
\item
  A-模 M 是生成模。
\item
  对每个 A-模 N，存在一个集合 I 和一个从 \(M^{(I)}\) 到 N 的满 A-线性映射。
\item
  存在一个整数 \(n \geqslant 0\) 和一个从 \(\mathbf{M}^n\) 到 \(\mathrm{A}_s\) 的满 A-线性映射。
\item
  存在整数 \(n \geqslant 0\)，使 \(A_s\) 同构于 \(M^n\) 的一个直因子子模。
\item
  迹理想 \(\tau (\mathbf{M})\) 等于 A。
\item
  存在整数 \(n \geqslant 0\)、M 的元素 \(x_{1}, \ldots, x_{n}\) 及 M* 的元素 \(x_{1}^{*}, \ldots, x_{n}^{*}\)，满足 \(\sum_{i=1}^{n} \langle x_{i}, x_{i}^{*} \rangle = 1\)。
\item
  \(\Rightarrow\) (ii)：设 M 是生成模，N 是一个左 A-模。存在一族从 M 到 N 的 A-线性映射 \((u_{i})_{i\in\mathrm{I}}\)，使我们有 \(N=\sum_{i\in\mathrm{I}}u_{i}(\mathrm{M})\)。映射 \((x_{i})\mapsto\sum_{i\in\mathrm{I}}u_{i}(x_{i})\)（从 \(\mathrm{M}^{(\mathrm{I})}\) 到 N）是 A-线性的满射。
\item
  \(\Rightarrow\) (iii)：设性质 (ii) 成立，取 \(N = A_s\)。于是我们有一个集合 I 和一个满线性映射 \(u: M^{(I)} \to A_s\)。由于 \(A_s\) 由元素 1 生成，存在 I 的有限子集 J 使 \(u(M^{(J)}) = A_s\)，从而得到 (iii)。
\item
  \(\Rightarrow\) (iv)：这由 II, §1, No.~12, 第 218 页的命题 21 得出。
\item
  \(\Rightarrow\) (v)：设 \(n \geqslant 1\) 是使 \(\mathrm{A}_s\) 同构于 M 的一个直因子子模的整数。我们有 \(\mathrm{A} = \tau(\mathrm{A}_s) \subset \tau(\mathrm{M}^n) = \tau(\mathrm{M})\)，因而 \(\tau(\mathrm{M}) = \mathrm{A}\)。
\item
  \(\Rightarrow\) (vi)：这是明显的。
\item
  \(\Rightarrow\) (i)：设 \(n\) 是整数 \(\geqslant 0\)，\(x_{1},\ldots ,x_{n}\) 是 M 的元素，\(x_{1}^{*},\ldots ,x_{n}^{*}\) 是 M \(^*\) 的元素，满足 \(\sum_{i = 1}^{n}\langle x_i,x_i^*\rangle = 1\)。设 N 是一个左 A-模，\(y\) 是 N 的一个元素。从 M 到 N 的映射 \(u_{i}:x\mapsto \langle x,x_{i}^{*}\rangle y\) 是 A-线性的，且我们有 \(y = \sum_{i = 1}^{n}u_{i}(x_{i})\)。这就证明 M 是一个生成 A-模。
\end{enumerate}

推论. —— 生成 A-模是忠实的。

设 \(a \in A\) 满足 \(aM = 0\)。利用定理 1 的蕴含关系 (i) \(\Rightarrow\) (iv)，我们得到 \(aA_s = 0\)，因此 \(a = 0\)。推论得证（II, §1, No.~12, 第 219 页）。

例. —— 1) A-模 \(A_s\) 是生成模。

\begin{enumerate}
\def\labelenumi{\arabic{enumi})}
\setcounter{enumi}{1}
\item
  每个非零自由 A-模都是生成模。更一般地，有生成商的模本身是生成模。
\item
  设 M 是一个反模为有限生成的半单 A-模。则 M 是生成 \(A_{M}\)-模。事实上，由 VIII, 第 8 页的引理 4，存在自然数 m 使 \((\mathrm{A}_{\mathrm{M}})_{s}\) 同构于 \(M^{m}\) 的一个子模。由于 \(M^{m}\) 是半单 \(A_{M}\)-模，\((\mathrm{A}_{\mathrm{M}})_{s}\) 同构于 \(M^{m}\) 的一个直因子子模，故 M 是生成 \(A_{M}\)-模（定理 1）。
\item
  设 A 是主理想整环，P 是有限生成 A-模。存在整数 \(n \geqslant 1\) 与 A 的递增理想序列 \((\mathfrak{a}_{i})_{1 \leqslant i \leqslant n}\)，使 P 同构于诸 \(A/\mathfrak{a}_{i}\) 的直和（VII, §4, No.~4, 第 19 页，定理 2）；P 的零化理想 a 等于 \(a_{1}\)。于是 P 是环 A/a 上的生成模。若 P 不是挠模，则我们有 a = 0，而 P 是生成 A-模。
\end{enumerate}

引理 1. —— 设 A 是交换环，M 是有限生成 A-模，Ann(M) 是它的零化子。设 \(\mathfrak{a}\) 是 A 的一个理想。下列性质等价：

\begin{enumerate}
\def\labelenumi{(\roman{enumi})}
\item
  aM = M.
\item
  Ann(M) + a = A.
\item
  存在一个元素 \(a\)（属于 \(\mathfrak{a}\)），使得 \(am = m\) 对每个 \(m \in \mathrm{M}\) 成立。 (i) \(\Rightarrow\) (ii)：设 \((x_1, \ldots, x_n)\) 是 A-模 M 的一个生成族。若 \(\mathfrak{aM} = \mathrm{M}\)，则每个 \(x_i\) 都可写成 \(\sum_{j=1}^{n} c_{ij} x_j\)，其中 \(c_{ij}\) 属于 \(\mathfrak{a}\)。把矩阵 \((c_{ij})\) 记作 \(C\)，把以 \(x_1, \ldots, x_n\) 为元素的列矩阵记作 \(X\)。我们有 \((I_n - C)X = 0\)。设 \(d\) 为行列式，\(V\) 为矩阵 \(I_n - C\) 的余子式矩阵。由 III, §8, No.~6, 第 532 页的公式 (26)，我们有 \(dX = {}^t V(I_n - C)X = 0\)，从而 \(d \in \operatorname{Ann}(\mathrm{M})\)。另一方面，由于 \(c_{ij}\) 属于 \(\mathfrak{a}\)，我们有 \(d \equiv 1 (\bmod \mathfrak{a})\)（III, §8, No.~5, 第 529 页，(18)），因而 \(1 \in \operatorname{Ann}(\mathrm{M}) + \mathfrak{a}\)。
\item
  \(\Rightarrow\) (iii)：设 (ii) 成立，则存在 \(a\in \mathfrak{a}\) 与 \(b\in \mathrm{Ann}(\mathrm{M})\) 使 \(a + b = 1\)。于是 \(am = m\) 对每个 \(m\in \mathrm{M}\) 成立。
\item
  \(\Rightarrow\) (i)：这是明显的。
\end{enumerate}

命题 3. —— 设 A 是交换环。每个有限生成的忠实投射 A-模都是生成模。更一般地，有限生成投射 A-模 P 是生成 \(A_{P}\)-模。

设 P 是有限生成投射 A-模。我们有 \(\tau(\mathrm{P})\mathrm{P}=\mathrm{P}\)（VIII, 第 80 页）。

若 A-模 P 是忠实的，则理想 \(\tau(\mathrm{P})\) 等于 A（引理 1），且 A-模 P 是生成模（定理 1）；这就得到第一条断言。第二条断言由下面的引理 2 得出。

引理 2. —— 设 A 是一个环，M 是投射 A-模。则 \(\mathrm{A_M}\)-模 M 是投射的。

设 \((x_{i})_{i\in\mathrm{I}}\) 是 A-模 M 的一个生成族。存在 A-模 M 上的一族线性型 \((x_{i}^{*})_{i\in\mathrm{I}}\)，使对每个 \(x\in\mathrm{M}\)，族 \((\langle x,x_{i}^{*}\rangle)_{i\in\mathrm{I}}\) 有有限支集，且 \(x=\sum_{i\in\mathrm{I}}\langle x,x_{i}^{*}\rangle x_{i}\)（II, §2, No.~6, 第 238 页，命题 12）。对每个 \(i\in\mathrm{I}\)，映射 \(x\mapsto\langle x,x_{i}^{*}\rangle_{\mathrm{M}}\) 是 \(A_{M}\)-模 M 上的线性型，且 \(x=\sum_{i\in\mathrm{I}}\langle x,x_{i}^{*}\rangle_{\mathrm{M}}x_{i}\) 对每个 \(x\in\mathrm{M}\) 成立。由同前所引，M 是投射 \(A_{M}\)-模。

